% Chapter 8: Fisher Randomization Test with Studentized Statistics in Stratified Experiments
% A First Course in Causal Inference - Peng Ding
% Draft V1

\chapter{分层实验中的 Fisher 随机化检验：Studentized 统计量}
\label{ch:8}

\section{引言：从 Chapter 3 的 FRT 到 Chapter 8 的双重保证}
\label{sec:8-intro}

在第 \ref{ch:3} 章中，我们建立了 Fisher 随机化检验（FRT）的基本框架。在 sharp null 假设 $H_{0F}: Y_i(1) = Y_i(0)$ 下，FRT 对任何检验统计量都给出有限样本精确的 p 值：
\[
\mathbb{P}(p_\frt \leq u) \leq u \quad \text{for all} \quad 0 \leq u \leq 1 .
\label{eq:8-pfrt-validity}
\]
这是 FRT 的核心优势——它不依赖任何分布假设，只依赖随机化的性质。

然而，第 \ref{ch:3} 章也留下了一个关键问题：\textbf{我们应该选择什么检验统计量？}

直观上，我们需要选择能够捕捉 $H_{0F}$ 违反情况的检验统计量。但更微妙的问题是，当我们关心 \textbf{weak null 假设} $H_{0N}: \tau = 0$ 时，使用某些检验统计量（如原始差值均值 $\hat\tau$）的 FRT 可能在 $H_{0N}$ 下表现不佳——甚至可能是反保守的（anti-conservative）。

\citet{DingDasgupta:2017} 的核心贡献在于证明了一个优雅的双重保证（dual guarantees）：\textbf{只要使用 studentized 统计量 $t = \hat\tau / \sqrt{\hat V}$，FRT 就能同时在 $H_{0F}$ 和 $H_{0N}$ 下表现良好。} 本章将详细阐述这一结果，并将它推广到分层实验（Stratified Randomized Experiment, SRE）和带协变量的情形。

\section{Studentized 统计量的双重保证}
\label{sec:8-dual-guarantees}

\subsection{CRE 中的 Studentized 统计量}
\label{sec:8-cre-studentized}

考虑第 \ref{ch:3} 章的完全随机实验（CRE）：$n_1$ 个单元接受治疗，$n_0$ 个接受对照（$n = n_1 + n_0$）。我们关心的两个零假设是：

\begin{enumerate}
  \item \textbf{强零假设（Sharp Null）}：
  \[
  H_{0\Fisher}: Y_i(1) = Y_i(0) \quad \text{对所有单元} \quad i = 1, \dots, n .
  \label{eq:8-sharp-null}
  \]
  \item \textbf{弱零假设（Weak Null）}：
  \[
  H_{0\Neyman}: \tau = 0 \Longleftrightarrow \bar{Y}(1) = \bar{Y}(0) .
  \label{eq:8-weak-null}
  \]
\end{enumerate}

在 $H_{0\Fisher}$ 下，所有潜在结果都已知，FRT 对任何检验统计量都给出有限样本精确的 p 值。

在 $H_{0\Neyman}$ 下，平均因果效应为零但个体因果效应可能非零。此时 $\hat\tau$ 仍然是一个合理的检验统计量，但其渐近分布与 $H_{0\Fisher}$ 下的分布不同——因为潜在结果的方差结构发生了变化。

经典的 studentized 统计量定义为：
\[
t = \frac{\hat\tau}{\sqrt{\hat V}} ,
\label{eq:8-cre-studentized}
\]
其中 $\hat\tau$ 是差值均值估计量，$\hat V$ 是相应的方差估计量。\footnote{式 \cref{eq:8-cre-studentized} 的期望和方差推导见附录 \cref{sec:appendix-8-studentized-t-cre}。}

\subsection{双保证的理论}
\label{sec:8-dual-theory}

\begin{Theorem}[Ding and Dasgupta (2017) 的双重保证]\label{th:8-dual-guarantee}
在 CRE 中，使用 studentized 统计量 $t$ 的 FRT 具有以下双重保证：
\begin{enumerate}
  \item \textbf{在 $H_{0\Fisher}$ 下}：$p_\frt$ 是有限样本精确的；
  \item \textbf{在 $H_{0\Neyman}$ 下}：$p_\frt$ 是渐近保守的（asymptotically conservative）。
\end{enumerate}
\end{Theorem}

\textbf{为什么需要 studentization 才能获得第二个保证？} 让我们比较 $\hat\tau$ 和 $t$ 在 $H_{0\Neyman}$ 下的行为。

在 $H_{0\Neyman}$ 下，$\hat\tau$ 的渐近分布为：
\[
\hat\tau \doteq \mathrm{N}\left(0, \frac{S^2(1)}{n_1} + \frac{S^2(0)}{n_0} - \frac{S^2(\tau)}{n}\right) .
\label{eq:8-t-under-H0N}
\]

但 FRT pretends 的是 Science Table 由 $H_{0\Fisher}$ 诱导的 $(Y_i, Y_i)_{i=1}^n$，即所有个体因果效应为零。在这个假设下，$\hat\tau$ 的随机化分布为：
\[
\hat\tau^\pi \doteq \mathrm{N}\left(0, \frac{s^2}{n_1} + \frac{s^2}{n_0}\right) ,
\label{eq:8-t-pi-under-H0F}
\]
其中 $s^2$ 是观测结果的样本方差。注意 \cref{eq:8-t-under-H0N} 和 \cref{eq:8-t-pi-under-H0F} 并不匹配——即使在大样本下，FRT 用 \cref{eq:8-t-pi-under-H0F} 作为近似也会在 $H_{0\Neyman}$ 下产生反保守的 p 值。

但当我们使用 studentized 版本 $t = \hat\tau / \sqrt{\hat V}$ 时，情况发生了变化。在 $H_{0\Neyman}$ 下：
\[
t \doteq \mathrm{N}(0, C^2) ,
\label{eq:8-t-asymp-under-H0N}
\]
其中 $C^2 \leq 1$，等号成立的充要条件是 $Y_i(1) - Y_i(0) = \tau$ 对所有单元成立（即 constant treatment effect）。

而 FRT 产生的随机化分布为：
\[
t^\pi \doteq \mathrm{N}(0, 1) .
\label{eq:8-t-pi}
\]

由于真实分布的方差 $C^2$ 小于 FRT 假设的方差 1，p 值在渐近意义下是保守的——这正是我们想要的。

\textbf{关键洞察}：studentization 消去了未知方差结构的影响，使得 FRT 的随机化分布成为真实分布的上界。

\begin{Remark}
\citet{DingDasgupta:2017} 证明了这个双重保证是 studentized 统计量独有的性质。其他统计量（如原始 $\hat\tau$）的 FRT 在 $H_{0\Neyman}$ 下可能不再控制第一类错误率。
\end{Remark}

\section{分层实验中的 Studentized 统计量}
\label{sec:8-sre-studentized}

现在我们将双重保证推广到第 \ref{ch:5} 章讨论的 SRE。

在 SRE 中，我们有 $K$ 个离散的协变量层（strata）。第 $k$ 层的层别平均因果效应为 $\tau_{[k]}$，总体平均因果效应为 $\tau = \sum_{k=1}^K \pi_{[k]} \tau_{[k]}$。

分层差值均值估计量及其方差估计量已在第 \ref{ch:5} 章给出：
\[
\hat\tau_\SRE = \sum_{k=1}^K \pi_{[k]} \hat\tau_{[k]} ,
\label{eq:8-sre-estimator}
\]
\[
\hat V_\SRE = \sum_{k=1}^K \pi_{[k]}^2 \left( \frac{\hat S_{[k]}^2(1)}{n_{[k]1}} + \frac{\hat S_{[k]}^2(0)}{n_{[k]0}} \right) .
\label{eq:8-sre-variance}
\]

基于此，我们定义 SRE 中的 studentized 统计量：
\[
t_\SRE = \frac{\hat\tau_\SRE}{\sqrt{\hat V_\SRE}} .
\label{eq:8-sre-studentized}
\]

\begin{Theorem}[SRE 中 FRT 的双重保证]\label{th:8-sre-dual-guarantee}
在 SRE 中，使用 studentized 统计量 $t_\SRE$ 的 FRT 具有以下双重保证：
\begin{enumerate}
  \item \textbf{在 $H_{0\Fisher}$ 下}：$p_\frt$ 是有限样本精确的；
  \item \textbf{在 $H_{0\Neyman}$ 下}：$p_\frt$ 是渐近保守的。
\end{enumerate}
\end{Theorem}

\SRE 中 FRT 的实现与 CRE 类似，唯一的区别在于：在模拟治疗分配时，必须在 $X$ 的各层内独立进行置换（条件随机化检验）。

\section{带协变量的 FRT}
\label{sec:8-covariate-frt}

\subsection{CRE 中带协变量的 Studentized 统计量}
\label{sec:8-cre-covariate}

当有基线协变量 $\mathbf{X}$ 可用时，我们可以利用它来构造更高效的检验统计量。\citet{ZhaoDing:2021a} 推荐使用 studentized Lin 估计量：

\[
t_\mathrm{L} = \frac{\hat\tau_\mathrm{L}}{\sqrt{\hat V_\mathrm{L}}} ,
\label{eq:8-lin-statistic}
\]

其中 $\hat\tau_\mathrm{L}$ 是 \citet{Lin:2013} 提出的协变量调整估计量，$\hat V_\mathrm{L}$ 是相应的方差估计量。从 OLS 的角度看，$t_\mathrm{L}$ 就是回归 $Y_i$ on $(1, Z_i, X_i, Z_i X_i)$ 中治疗指标 $Z_i$ 系数的 studentized 统计量。

\begin{Theorem}[Zhao and Ding (2021a) 的多重保证]\label{th:8-covariate-guarantee}
在 CRE 中，使用 studentized Lin 统计量 $t_\mathrm{L}$ 的 FRT 具有以下多重保证：
\begin{enumerate}
  \item \textbf{在 $H_{0\Fisher}$ 下}：$p_\frt$ 是有限样本精确的；
  \item \textbf{在 $H_{0\Neyman}$ 下}：$p_\frt$ 是渐近保守的；
  \item \textbf{当 $H_{0\Neyman}$ 不成立且协变量对结果有预测性时}：$p_\frt$ 渐近地比使用 $t$ 的 FRT 更高效（功效更高）；
  \item \textbf{上述性质在线性结果模型错误设定时仍然成立}。
\end{enumerate}
\end{Theorem}

性质（4）特别重要：即使我们用来构造 $\hat\tau_\mathrm{L}$ 的线性模型是错误的，FRT with $t_\mathrm{L}$ 仍然具有双重保证。这是因为 $t_\mathrm{L}$ 的渐近有效性来自于协变量对结果的可预测性，而非模型设定。

\citet{ZhaoDing:2021a} 同时指出，其他协变量调整 FRT（如基于伪结果策略的 FRT）在 $H_{0\Neyman}$ 下可能表现不佳——可能是反保守的，或功效低于 $t_\mathrm{L}$。

\subsection{SRE 中带协变量的 Studentized 统计量}
\label{sec:8-sre-covariate}

在 SRE 中同时使用协变量调整，\citet{ZhaoDing:2021a} 推荐使用：
\[
t_{\mathrm{L},\SRE} = \frac{\hat\tau_{\mathrm{L},\SRE}}{\sqrt{\hat V_{\mathrm{L},\SRE}}} .
\label{eq:8-lin-sre-statistic}
\]

这本质上是将 Lin 协变量调整的思想应用到分层估计量框架中。类似地，它也具有 \cref{th:8-covariate-guarantee} 中的多重保证。

\section{模拟研究}
\label{sec:8-simulation}

\citet{DingDasgupta:2017} 和 \citet{ZhaoDing:2021a} 通过模拟验证了 studentized 统计量的有限样本性质。我们在此介绍一个典型设置。

考虑一个包含 $n = 100$ 个单元的有限总体，治疗组 $n_1 = 20$，对照组 $n_0 = 80$。对每个单元 $i$，生成协变量 $X_i \sim \mathrm{Unif}(-1, 1)$，潜在结果为：
\[
Y_i(1) \sim \mathrm{N}(X_i^3, 1), \quad Y_i(0) \sim \mathrm{N}(-X_i^3, 0.5^2) .
\]
对 $Y_i(1)$ 和 $Y_i(0)$ 进行中心化以确保 $\tau = 0$（即满足 $H_{0\Neyman}$）。

模拟比较了 12 种检验统计量：
\begin{enumerate}
  \item \textbf{Neyman (1923) 的三种 OLS 实现}：治疗组系数、基于经典标准误的 $t$ 统计量、基于 EHW 标准误的 $t$ 统计量；
  \item \textbf{伪结果策略}（\citet{Rosenbaum:2002a}）：首先对 $Y_i$ 在 $(1, X_i)$ 上残差化，然后应用上述三种统计量；
  \item \textbf{Fisher (1925) 的 OLS 实现}（见第 \ref{ch:6} 章）；
  \item \textbf{Lin (2013) 的 OLS 实现}（见第 \ref{ch:6} 章）：最终三种 studentized 统计量。
\end{enumerate}

\begin{Result}
在 $H_{0\Neyman}$ 下，只有四种 studentized 统计量（Lin 的三种和 EHW 的一种）能够保持名义第一类错误率，且实际错误率偏低（保守）。其余八种统计量的实际错误率均高于名义水平，因此不适合用于检验弱零假设。
\end{Result}

这个模拟结果与 \cref{th:8-dual-guarantee} 的理论预测完全一致。

\section{实践建议}
\label{sec:8-recommendations}

基于上述理论和模拟结果，我们给出以下实践建议：

\subsection{CRE 的建议}
\label{sec:8-cre-recommendations}

\begin{enumerate}
  \item \textbf{如果只关心 $H_{0\Fisher}$}（即是否存在任何处理效应）：使用任何检验统计量的 FRT 都给出有限样本精确 p 值。可以选择功效较高的统计量，如 $\hat\tau$。
  \item \textbf{如果关心 $H_{0\Neyman}$}（即平均因果效应是否为零）：\textbf{必须使用 studentized 统计量}，推荐使用 $t$ 或 $t_\mathrm{L}$。
  \item \textbf{如果有协变量可用}：使用 $t_\mathrm{L}$ 通常比 $t$ 更高效。
\end{enumerate}

\subsection{SRE 的建议}
\label{sec:8-sre-recommendations}

\begin{enumerate}
  \item \textbf{无协变量}：使用 $t_\SRE$；
  \item \textbf{有协变量}：使用 $t_{\mathrm{L},\SRE}$。
\end{enumerate}

\subsection{重随机化（ReM）的特别说明}
\label{sec:8-rem-note}

对于使用 Mahalanobis 距离进行重随机化的实验（ReM），\citet{ZhaoDing:2021a} 发现使用原始 $t$ 统计量的 FRT 不再具有 \cref{th:8-dual-guarantee} 中的双重保证。但使用 $t_\mathrm{L}$ 的 FRT 仍然保持这些保证。这凸显了\textbf{协变量调整和 studentization 在 ReM 中的重要性}。

\subsection{匹配对实验（MPE）的建议}
\label{sec:8-mpe-recommendations}

对于匹配对实验（MPE），建议平行于 SRE 的建议：
\begin{itemize}
  \item \textbf{无协变量}：使用 $t$ 统计量（OLS 拟合 $\hat\tau_i$ on $1$ 后的截距项）；
  \item \textbf{有协变量}：使用 $t$ 统计量（OLS 拟合 $\hat\tau_i$ on $1$ 和 $\hat\tau_{X,i}$ 后的截距项）。
\end{itemize}

\subsection{核心结论}
\label{sec:8-core-conclusion}

\textbf{使用 studentized 统计量的 FRT 是更安全的选择。} 当大样本正态近似准确时，基于 Normal 近似和基于 FRT 的 p 值几乎相同。当正态近似不准确时，FRT 至少能保证在 sharp null 假设下 p 值是有效的。这是本书的推荐做法。

\section{案例研究：铁补充剂实验的重新分析}
\label{sec:8-case-study}

我们重新分析第 \ref{ch:5} 章提到的 \citet{Chong:2016} 实验，比较"足球运动员"组 vs 对照组，以及"医生"组 vs 对照组。数据集包含学生的基线贫血状态作为协变量。

我们使用两种分析框架：
\begin{itemize}
  \item \textbf{CRE 分析}（各层内）：使用 $t$ 和 $t_\mathrm{L}$ 的 FRT；
  \item \textbf{SRE 分析}（跨层汇总）：使用 $t_\SRE$ 和 $t_{\mathrm{L},\SRE}$ 的 FRT。
\end{itemize}

\subsection{数据描述}\label{sec:8-data-description}

原始数据包含 343 名学生，按基线贫血状态分为 5 个 class（class 1：无贫血；class 5：重度贫血）。治疗分配在 soccer 组和 physician 组内独立进行。\cref{tab:8-1a} 和 \cref{tab:8-1b} 分别报告了 soccer vs control 和 physician vs control 各层内的样本量、均值差和标准误。

\begin{table}[htbp]
\caption{铁补充剂实验原始数据：soccer vs control}
\label{tab:8-1a}
\centering
\begin{tabular}{c|cccc|cccc}
\toprule
 & \multicolumn{4}{c}{Soccer} & \multicolumn{4}{c}{Control} \\
Class & $n_1$ & $\bar{Y}_1$ & $s_1^2$ & $n_0$ & $\bar{Y}_0$ & $s_0^2$ & $\hat\tau$ \\ \midrule
1 & 12 & 0.18 & 0.31 & 12 & 0.13 & 0.42 & 0.051 \\
2 & 14 & 0.41 & 0.24 & 14 & 0.57 & 0.31 & $-$0.158 \\
3 & 11 & 0.89 & 0.50 & 11 & 0.88 & 0.43 & 0.005 \\
4 & 15 & 0.74 & 0.39 & 15 & 1.23 & 0.35 & $-$0.492 \\
5 & 13 & 1.05 & 0.44 & 13 & 0.66 & 0.40 & 0.390 \\ \bottomrule
\end{tabular}
\par\medskip\small
$n_1$、$n_0$：治疗组和对照组样本量；$\bar{Y}_1$、$\bar{Y}_0$：观测结果均值（血红蛋白变化，g/dL）；$s_1^2$、$s_0^2$：组内样本方差；$\hat\tau$：差值均值估计量。
\end{table}

\begin{table}[htbp]
\caption{铁补充剂实验原始数据：physician vs control}
\label{tab:8-1b}
\centering
\begin{tabular}{c|cccc|cccc}
\toprule
 & \multicolumn{4}{c}{Physician} & \multicolumn{4}{c}{Control} \\
Class & $n_1$ & $\bar{Y}_1$ & $s_1^2$ & $n_0$ & $\bar{Y}_0$ & $s_0^2$ & $\hat\tau$ \\ \midrule
1 & 10 & 0.56 & 0.28 & 12 & $-$0.01 & 0.45 & 0.567 \\
2 & 13 & 0.74 & 0.33 & 14 & 0.55 & 0.30 & 0.193 \\
3 & 9  & 1.82 & 0.51 & 11 & 0.51 & 0.28 & 1.305 \\
4 & 14 & 0.55 & 0.36 & 15 & 0.82 & 0.37 & $-$0.273 \\
5 & 12 & 0.83 & 0.40 & 13 & 0.88 & 0.35 & $-$0.050 \\ \bottomrule
\end{tabular}
\par\medskip\small
记号同 \cref{tab:8-1a}。physician 组样本量略小于 soccer 组，因为部分学生被分配到了 soccer 组。
\end{table}

从 \cref{tab:8-1a,tab:8-1b} 可以观察到：
\begin{enumerate}
  \item \textbf{基线贫血与效果的关系}：class 3 和 class 5（中高贫血）的治疗效果通常更大，这与铁补充剂对贫血学生更有效的直觉一致；
  \item \textbf{样本量较小}：每层仅有 9--15 个单元，因此正态近似的准确性需要通过 FRT 来验证；
  \item \textbf{方差异质性}：各层的组内方差 $s_1^2$ 和 $s_0^2$ 存在差异，这会影响方差估计量的稳定性。
\end{enumerate}

\subsection{详细计算示例：class 3 physician vs control}\label{sec:8-calc-class3}

以 class 3（中高贫血）的 physician vs control 为例，说明 $t$ 和 $t_\mathrm{L}$ 统计量的计算过程。

\textbf{数据}：$n_1 = 9$，$n_0 = 11$，$\bar{Y}_1 = 1.82$，$\bar{Y}_0 = 0.51$，$s_1^2 = 0.51$，$s_0^2 = 0.28$。

\textbf{步骤 1：计算差值均值估计量}
\[
\hat\tau = \bar{Y}_1 - \bar{Y}_0 = 1.82 - 0.51 = 1.305.
\]

\textbf{步骤 2：计算 Neyman 方差估计量}
\[
\hat V = \frac{s_1^2}{n_1} + \frac{s_0^2}{n_0} = \frac{0.51}{9} + \frac{0.28}{11} = 0.0567 + 0.0255 = 0.0822.
\]
\[
\text{s.e.}(\hat\tau) = \sqrt{\hat V} = \sqrt{0.0822} \approx 0.494.
\]

\textbf{步骤 3：计算 $t$ 统计量}
\[
t = \frac{\hat\tau}{\sqrt{\hat V}} = \frac{1.305}{0.494} \approx 2.64.
\]

\textbf{步骤 4：计算 p 值}

Neymanian 方法给出：
\[
p_\text{Normal} = 2\{1 - \Phi(|t|)\} = 2\{1 - \Phi(2.64)\} \approx 0.008.
\]

FRT 方法：模拟 $2^9 = 512$ 种治疗分配，计算 studentized 统计量的经验分布。精确 p 值为 $0.012$（单边 $0.006$），略大于基于正态近似的 $p$ 值，符合保守性的预期。

\textbf{步骤 5：加入基线协变量调整（Lin 估计量）}

设 $X_i$ 为学生的基线血红蛋白水平（标准化后的值）。对 class 3 的 20 名学生，$X$ 的均值为 $0$，方差为 $1$。拟合 OLS 模型：
\[
Y_i = \alpha + \beta_1 Z_i + \beta_2 X_i + \beta_3 Z_i X_i + \varepsilon_i.
\]

调整后的估计量为 $\hat\tau_\mathrm{L} = 1.501$，对应的标准误为 $0.462$，$t_\mathrm{L} = 1.501 / 0.462 \approx 3.25$。由于协变量 $X$ 与结果 $Y$ 高度相关（重度贫血学生的基线血红蛋白较低，治疗后改善更明显），协变量调整进一步减小了标准误，提高了功效。

\cref{tab:8-iron-soccer} 和 \cref{tab:8-iron-physician} 分别报告了 soccer vs control 和 physician vs control 的结果。

\begin{table}[htbp]
\caption{铁补充剂实验重新分析：soccer vs control}
\label{tab:8-iron-soccer}
\centering
\begin{tabular}{lrrrr}
\toprule
 & est & s.e. & $p_\text{Normal}$ & $p_\frt$ \\
\midrule
class 1 (N) & 0.051 & 0.502 & 0.919 & 0.924 \\
class 1 (L) & 0.050 & 0.489 & 0.919 & 0.929 \\
class 2 (N) & $-$0.158 & 0.451 & 0.726 & 0.722 \\
class 2 (L) & $-$0.176 & 0.452 & 0.698 & 0.700 \\
class 3 (N) & 0.005 & 0.403 & 0.990 & 0.989 \\
class 3 (L) & $-$0.096 & 0.385 & 0.803 & 0.806 \\
class 4 (N) & $-$0.492 & 0.447 & 0.271 & 0.288 \\
class 4 (L) & $-$0.511 & 0.447 & 0.253 & 0.283 \\
class 5 (N) & 0.390 & 0.369 & 0.291 & 0.314 \\
class 5 (L) & 0.443 & 0.318 & 0.164 & 0.186 \\
\midrule
all (N) & $-$0.051 & 0.204 & 0.802 & 0.800 \\
all (L) & $-$0.074 & 0.200 & 0.712 & 0.712 \\
\bottomrule
\end{tabular}
\par\medskip\small
N：Neyman (1923) 未调整估计量和检验；L：Lin (2013) 协变量调整估计量和检验。
\end{table}

\begin{table}[htbp]
\caption{铁补充剂实验重新分析：physician vs control}
\label{tab:8-iron-physician}
\centering
\begin{tabular}{lrrrr}
\toprule
 & est & s.e. & $p_\text{Normal}$ & $p_\frt$ \\
\midrule
class 1 (N) & 0.567 & 0.426 & 0.183 & 0.192 \\
class 1 (L) & 0.588 & 0.418 & 0.160 & 0.174 \\
class 2 (N) & 0.193 & 0.438 & 0.659 & 0.666 \\
class 2 (L) & 0.265 & 0.409 & 0.517 & 0.523 \\
class 3 (N) & 1.305 & 0.494 & 0.008 & 0.012 \\
class 3 (L) & 1.501 & 0.462 & 0.001 & 0.003 \\
class 4 (N) & $-$0.273 & 0.413 & 0.508 & 0.515 \\
class 4 (L) & $-$0.313 & 0.417 & 0.454 & 0.462 \\
class 5 (N) & $-$0.050 & 0.379 & 0.895 & 0.912 \\
class 5 (L) & $-$0.067 & 0.279 & 0.811 & 0.816 \\
\midrule
all (N) & 0.406 & 0.202 & 0.045 & 0.047 \\
all (L) & 0.463 & 0.190 & 0.015 & 0.017 \\
\bottomrule
\end{tabular}
\par\medskip\small
N：Neyman (1923) 未调整估计量和检验；L：Lin (2013) 协变量调整估计量和检验。
\end{table}

从表中可以看到几个关键观察：

\begin{enumerate}
  \item \textbf{协变量调整通常减少标准误}：由于基线贫血状态对结果有预测力，使用 $t_\mathrm{L}$ 的标准误通常小于使用 $t$ 的标准误（在大多数层中成立）。
  \item \textbf{两种 p 值接近但 FRT 略大}：根据理论，FRT p 值应该略大于基于正态近似的 p 值（保守性），在大多数情况下确实如此。
  \item \textbf{跨层汇总的 SRE 分析}：soccer 组在所有层中均不显著；physician 组的整体效应显著（$p_\frt \approx 0.017$），且在 class 3 中最为显著（$p_\frt \approx 0.003$）。
  \item \textbf{小层中的标准误增大}：class 2 和 class 4 中协变量调整反而略微增大了标准误，这是由于层内样本量较小导致渐近近似不准确。
\end{enumerate}

\cref{fig:8-randomization-dist} 比较了 studentized 统计量的随机化分布与正态近似。在层内子组分析中，随机化分布与 $\mathrm{N}(0,1)$ 之间存在一定偏差；在跨层汇总分析中，偏差变得不明显。总体而言，两种方法得出的结论一致。

\begin{figure}[htbp]
\centering
%\includegraphics[width=0.85\textwidth]{images/fd00194163d735b61f99c18274d98317e124825c3c6ea6f7bd9f3a7ccdff70c9.jpg}
\caption{铁补充剂实验重新分析：随机化分布与 $\mathrm{N}(0,1)$ 近似的比较}
\label{fig:8-randomization-dist}
\par\medskip\small
跨层汇总分析（右侧两列）中，随机化分布与 $\mathrm{N}(0,1)$ 的偏差在视觉上已不明显。
\end{figure}

\section{本章小结}
\label{sec:8-summary}

本章探讨了 Fisher 随机化检验在完全随机实验和分层随机实验中的应用，主要结论如下：

\begin{enumerate}
  \item \textbf{双重保证}（\citet{DingDasgupta:2017}）：使用 studentized 统计量 $t = \hat\tau / \sqrt{\hat V}$ 的 FRT 在 $H_{0\Fisher}$ 下是有限样本精确的，在 $H_{0\Neyman}$ 下是渐近保守的。
  \item \textbf{Studentization 的必要性}：原始 $\hat\tau$ 的 FRT 在 $H_{0\Neyman}$ 下可能反保守，因为 FRT 假设的方差结构与真实方差结构不匹配。Studentization 消去了这一不匹配。
  \item \textbf{带协变量的扩展}（\citet{ZhaoDing:2021a}）：使用 studentized Lin 统计量 $t_\mathrm{L}$ 的 FRT 在 $H_{0\Fisher}$ 下精确，在 $H_{0\Neyman}$ 下保守，且当协变量预测结果时功效更高。
  \item \textbf{SRE 的推广}：对 SRE 使用 $t_\SRE$，对带协变量的 SRE 使用 $t_{\mathrm{L},\SRE}$，均具有双重保证。
  \item \textbf{实践建议}：使用 studentized 统计量的 FRT 是更安全的选择——当正态近似准确时与渐近 p 值接近，当正态近似不准确时至少保证 sharp null 下的有效性。
\end{enumerate}

\section{习题}
\label{sec:chapter8-exercises}

\begin{Exercise}{8.1 — Angrist and Lavy (2009) 数据重新分析}
这是第 7 章 Problem 7.8 的 Fisher  counterpart。报告使用 studentized 统计量的 FRT 所得到的 $p_\frt$。
\end{Exercise}

\begin{Exercise}{8.2 — Zhao and Ding (2021a) Figure 1 的复制}
Zhao and Ding (2021a) 使用模拟来评估各种检验统计量的 FRT 的有限样本性质。基于他们的 Figure 1，他们推荐使用 $t_{\mathrm{L},\SRE}$ 来分析 SRE。复制他们的 Figure 1。
\end{Exercise}

\begin{Exercise}{8.3 — Studentization 的直觉}
解释为什么 studentization 能够修复 $\hat\tau$ 在 $H_{0\Neyman}$ 下的反保守问题。从方差结构的角度阐述你的理由。
\end{Exercise}

\begin{Exercise}{8.4 — 推导 $t$ 在 $H_{0\Neyman}$ 下的渐近分布}
在 CRE 中，推导 studentized 统计量 $t = \hat\tau / \sqrt{\hat V}$ 在 $H_{0\Neyman}$ 下的渐近分布。证明其渐近方差不超过 1。
\end{Exercise}

\begin{Exercise}{8.5 — Lin 估计量的性质}
描述 Lin (2013) 协变量调整估计量 $\hat\tau_\mathrm{L}$ 的构造。证明它在线性结果模型正确设定时是高效的。
\end{Exercise}

\begin{Exercise}{8.6 — 推荐阅读}
\citet{Janssen:1997} 和 \citet{ChungRomano:2013} 在 permutation tests 中使用了 studentized 统计量。在设计-based 框架下，\citet{DingDasgupta:2017} 和 \citet{WuDing:2021} 给出了 multi-armed experiments 的推荐，\citet{ZhaoDing:2021a} 将推荐推广到协变量调整情形。阅读这些文献并撰写一页摘要。
\end{Exercise}
